\section{Capítulo~\ref{sec:recognition}. Reconocimiento}

La velocidad del reconocimiento, la estructura de las primeras etapas, la lectura lineal de la respuesta y el aprendizaje a partir de la continuidad del tiempo están medidos en humanos y en monos; la reducción de toda la cadena a dos operaciones y el aprendizaje a partir de video se comprobaron en los modelos del libro, y las explicaciones de las ilusiones van de bien fundadas a discutidas.

{\footnotesize
\setlength{\tabcolsep}{3pt}\renewcommand{\arraystretch}{1.15}
\begin{longtable}{>{\raggedright}p{0.5cm}>{\raggedright}p{4.0cm}>{\raggedright}p{5.6cm}>{\raggedright}p{2.3cm}>{\raggedright\arraybackslash}p{1.7cm}}
\toprule
N.º & Proposición & Experimento y qué se midió & Fuente & Estado \\
\midrule\endfirsthead
\toprule
N.º & Proposición & Experimento y qué se midió & Fuente & Estado \\
\midrule\endhead
1 & Con traslaciones, la comparación píxel a píxel con un patrón no acierta más que una moneda; el par de etapas~\eqref{eq:scstage} reconoce la figura en cualquier lugar & \texttt{models/\allowbreak recognition\_models.py}, “retina” de $24$ puntos, $4000$ pruebas: patrón, $0.49$, $0.51$, $0.52$ con fracciones de puntos invertidos $0$, $0.1$, $0.2$; par $S$ y $C$, $1.00$, $0.86$, $0.70$ & deducción en el capítulo & modelo \\
2 & Un animal en una imagen se reconoce en $\sim150$\,ms, una pasada por una decena de etapas de $10$--$20$\,ms & Humanos, tarea “hay un animal o no” con fotografías mostradas $20$\,ms: los potenciales evocados de las dos respuestas divergen a los $\sim150$\,ms. Mono: el tiempo de la primera respuesta crece de etapa en etapa (V1 antes que todas, luego V2 y V4). El número de etapas y “cero o una espiga por etapa” se deducen de estas cifras & \cite{Thorpe1996,Schmolesky1998} & medido \\
3 & La etapa $S$ es un AND de partes; la etapa $C$, el máximo sobre posiciones (proposición~\ref{prop:conveyor}) & Gato, corteza visual primaria: las células simples responden a un borde de cierta orientación en cierto lugar; las complejas, a la misma orientación en una zona mayor. Registro intracelular de células complejas: en muchas, la respuesta a dos barras se acerca a la mayor de las dos respuestas, en promedio lo más cerca de la operación max. El máximo mediante inhibición mutua es la proposición~\ref{prop:wta}; la división por la actividad total, una revisión & \cite{HubelWiesel1962,Lampl2004,Carandini2012} & medido \\
4 & Más arriba en la cadena, las combinaciones son más grandes y complejas, y la respuesta depende cada vez menos de la posición & Mono, simplificación de estímulos hasta el “rasgo crítico”: las células de V4 y de la corteza temporal inferior posterior responden a combinaciones complejas de forma, color y textura con campo pequeño; en la corteza temporal inferior anterior el campo es mayor. La alternancia de $S$ y $C$ en un modelo reproduce este cuadro & \cite{KobatakeTanaka1994,Riesenhuber1999} & medido \\
5 & Las redes convolucionales repiten esta alternancia y predicen mejor que nada las respuestas de las etapas superiores; el objeto se lee linealmente de las frecuencias de un grupo de neuronas & Una red entrenada para reconocer, sin ajuste al cerebro: la capa superior predice las respuestas de las neuronas de la corteza temporal inferior del mono; las intermedias, las respuestas de V4. Un clasificador lineal sobre la actividad de $\sim100$ células elegidas al azar, en ventanas desde $12.5$\,ms, reconoce el objeto y la categoría con distintas posiciones y tamaños & \cite{Fukushima1980,Yamins2014,Hung2005} & medido \\
6 & La regla de Hebb con traza~\eqref{eq:traceinvariance} enseña la invariancia a partir de video sin etiquetas si la traza no es más corta que $\sim20$\,ms & La idea de los rasgos lentos y de la regla con traza es teoría. \texttt{models/\allowbreak recognition\_models.py}, dos neuronas $C$, $16$ entradas, $10$ corridas, pausa entre objetos de $0.3$\,s. Con $\tau_{\text{tr}}=5$\,ms la coincidencia con la figura es de $0.52$ en promedio, con $10$\,ms, $0.56$; con $20$, $50$ y $200$\,ms, $1.00$ en todas las corridas (con $30$\,ms, una corrida de diez se equivoca) & \cite{Foldiak1991,WiskottSejnowski2002} & modelo \\
7 & En el cerebro la invariancia se aprende de la continuidad del tiempo & Mono: el objeto se sustituía sin que se notara durante un salto del ojo hacia un lugar determinado del campo; la invariancia de posición de las neuronas de la corteza temporal inferior cambiaba de acuerdo con la sustitución, de forma significativa ya al cabo de una hora. Que sea la regla principal del aprendizaje visual es una hipótesis & \cite{LiDiCarlo2008} & medido \\
8 & Movimiento: detectores de dirección y luego el mismo par AND/OR sobre zonas grandes & Detectores de dirección en la retina del conejo; en el área MT del mono, las $110$ neuronas registradas son selectivas a la dirección, y la respuesta al movimiento es más fuerte que en V1 & \cite{Barlow1965,Albright1984} & medido \\
9 & Las etapas superiores envían hacia abajo una predicción, y hacia arriba va el error & El modelo explica los efectos extraclásicos de los campos de la corteza visual primaria; algunas observaciones concuerdan (revisión); como estructura general de la cadena, no está demostrado & \cite{RaoBallard1999,Keller2018} & hipótesis \\
10 & Las imágenes difíciles requieren realimentación y varias vueltas & Mono: para las imágenes “difíciles”, la decisión en la corteza temporal inferior se forma $\sim30$\,ms más tarde; estas respuestas tardías las predice mal una red directa y mejor las redes con realimentación o más profundas & \cite{Kar2019} & medido \\
11 & Una alteración de píxeles invisible engaña a la red; la visión humana es mucho más robusta, pero no invulnerable & En redes: una perturbación pequeña, elegida a propósito, cambia la respuesta; el ejemplo “panda $\to$ gibón” es del segundo trabajo. En humanos, con presentación breve, las falsificaciones que se transfieren bien entre redes desvían la elección & \cite{Szegedy2014,Goodfellow2015,Elsayed2018} & medido \\
12 & Ilusiones de expectativa: la entrada poco fiable cede ante la expectativa; lo pálido se mueve “más despacio”, el vestido se ve de formas distintas (sección~\ref{sec:illusions}) & Las rejillas de menor contraste parecen moverse más despacio. Encuesta a $1401$ personas: el vestido se ve en tres variantes estables, y una pista sobre la iluminación cambia el color; en $\sim13\,000$ observadores decide la suposición sobre la iluminación. Un modelo bayesiano con expectativa de movimientos lentos explica varias ilusiones de movimiento & \cite{Thompson1982,LaferSousa2015,Wallisch2017,Weiss2002,Purves2011} & hipótesis \\
13 & Control de ganancia: la cascada, la inclinación y el contraste; la rejilla de Hermann no es inhibición de vecinos & Los detectores de dirección de la retina del conejo reducen su frecuencia ante un movimiento prolongado y, al detenerse, caen por debajo del fondo. La ilusión de inclinación se reproduce con un modelo de división por la actividad de los vecinos. Modificaciones de la rejilla que no cambian la respuesta de las neuronas de salida de la retina eliminan las manchas: se descarta la explicación por inhibición de vecinos & \cite{BarlowHill1963,Schwartz2009,SchillerCarvey2005} & medido \\
14 & Compensación de retardos: un objeto en movimiento parece ir por delante del destello & La ilusión está medida. La psicofísica contradice tanto la prolongación del movimiento hacia adelante como la diferencia de latencias; se propuso una explicación retrospectiva, según los sucesos de $\sim80$\,ms después del destello. El debate no está resuelto & \cite{Nijhawan1994,EaglemanSejnowski2000} & hipótesis \\
15 & Las “serpientes” inmóviles giran: los detectores de dirección leen la diferencia de latencias~\eqref{eq:latency} & La ilusión funciona también sin color; los pares de elementos vecinos la generan por separado; las neuronas corticales del mono selectivas a la dirección dan una respuesta direccional a estos pares inmóviles. En humanos, el giro lo disparan las microsacadas y los parpadeos & \cite{Conway2005,OteroMillan2012} & medido \\
16 & Imágenes ambiguas: inhibición mutua, fatiga y ruido; los cambios los provoca sobre todo el ruido & Modelo de grupos de neuronas en competencia: en él los cambios los provoca el \emph{ruido}, y sin ruido no ocurren; explica el aumento de la frecuencia de cambios con la fuerza del estímulo. La fatiga tiene aquí un papel menor & \cite{MorenoBote2007} & hipótesis \\
17 & Parte de las ilusiones es consecuencia de la tarea que aprende la máquina & Una red entrenada para predecir el cuadro siguiente de un video “ve” el giro de los anillos inmóviles y no lo ve en las imágenes de control; las redes para limpiar imágenes y aumentar su nitidez caen en ilusiones de color y brillo, como las personas & \cite{Watanabe2018,GomezVilla2020} & medido \\
\bottomrule
\end{longtable}}
